% ===================================================================
%  نقشہ نمبر 8 — واڈھی۔ — شکار، انگوراں دی واڈھی، مچھی پھڑنا۔
%  Traduction 2026 d'apres texto/io, controlee sur texto/fr, texto/en
%  et texto/hi. Consigne : voir texto/pnb/10-naqsha-01.tex.
%
%  L'appel de note est dans le TITRE de la scene — « واڈھی (*) » — et
%  non dans un alinea : c'est ainsi que l'ido le compose.
% ===================================================================

%%K t08-noto-1 noto
\VUnotes{143.2mm}{%
\textit{(*) کیوں جے ایہہ لفظ صاف نئیں : سݨ دی واڈھی؟ انگوراں دی؟ سیباں دی؟ شاید
چنگا ہووے پئی «} \VUgras{mesono} \textit{» کم وچ لیایا جاوے، جیہدی صلاح}
Mondo \textit{XIV، صفحہ 242 وچ دِتی گئی اے۔ پر ہُݨ تیکر اوس نویں دھات نوں
اکیڈمی نے اپݨایا نئیں تے اوہ ایدو لہر دے عام قاعدیاں مطابق بحث دی اُڈیک وچ اے۔}%
}

%%K t08-apar-1 apar
\VUsaut{5.0mm}
\VUtitre{15.0pt}{60}{نقشہ نمبر 8}
\VUsaut{3.0mm}
\VUcentre{13.2pt}{90}{\textit{پہلا منظر۔}}
\VUsaut{3.0mm}
\VUpk{10.2pt}{40}{واڈھی (*)۔}
\VUsaut{4.0mm}
\VUpk{10.2pt}{40}{پنڈ دا رُوپ۔}
\VUsaut{3.0mm}

%%K t08-c1-01-1 p
1. --- \VUgras{گرمی} نال پنڈ نے فیر اک واری اپݨا رُوپ بدل لیا اے۔ سیاڑ نوں سونپیا
بی پھُٹ پیا؛ زمین وچوں اک نِکا ہرا بُوٹا نکلیا تے اوہنے اپݨا سِٹا دِتا۔ داݨا
بݨیا، فیر پکیا، تے ہُݨ واڈھی توں بغیر کجھ باقی نئیں۔

%%K t08-c1-02-1 p
2. --- \VUgras{جنگلی پھُل} کھِڑ آئے نیں۔ \VUgras{نیل پھُل} \textsuperscript{(1)}،
\VUgras{پوست دے پھُل} \textsuperscript{(2)} تے \VUgras{ڈِجیٹیلس}
\textsuperscript{(3)} کݨک دیاں پیلیاں دے اکو جیہے رنگ اُتے اپݨے تازے
\VUgras{پھُلاں} دا خوش وِرودھ رکھدے نیں۔

%%K t08-c1-03-1 p
3. --- پھلاں والے رُکھاں نے پتّر کڈ لئے نیں؛ پھل پک گئے نیں۔ نِکا
\VUgras{چیری دا رُکھ} \textsuperscript{(4)} سوہݨیاں \VUgras{چیریاں}
\textsuperscript{(5)} نال لدیا اے، جیہناں نوں نِکے توڑدے تے چاء نال کھاندے نیں۔

%%K t08-c1-04-1 p
4. --- اِجّڑ ہُݨ سارا دن باہر لنگھا سکدا اے۔

%%K t08-c1-04-2 p
\VUgras{لیلے} \textsuperscript{(6)} وڈے ہو کے سوہݨیاں \VUgras{بھیڈاں}
\textsuperscript{(7)} تے سوہݨے \VUgras{چھترے} \textsuperscript{(8)} بݨ گئے نیں، سنی
اُون والے۔ اوہ \VUgras{چراگاہ} \textsuperscript{(9)} دا \VUgras{گھاہ} شانتی نال
چردے نیں، جیدے اُتے \VUgras{گُل بہار} کھِلرے نیں۔

%%K t08-c1-04-3 p
چھیتی ای ایہناں ڈنگراں دی \VUgras{اُون} لاہی جاوے گی، جیدے توں ساڈے کپڑیاں دا
کپڑا بݨدا اے۔

%%K t08-tit-2 sub
\VUsaut{5.0mm}
\VUpk{10.2pt}{40}{پالی۔}
\VUsaut{3.4mm}

%%K t08-c1-05-1 p
5. --- \VUgras{پالی} \textsuperscript{(10)} اوہناں ول ودھ دھیان دیندا نئیں لگدا۔
اوہدی اکو فکر اے اُفق ول جاندا طوفان، تے \VUgras{آجڑی} \textsuperscript{(13)}،
جیہڑا اپݨی \VUgras{صراحی} \textsuperscript{(14)} بُلھاں تیکر چُکی ہویا اے، اوہدے
توں وی ودھ بے فکر لگدا اے۔

%%K t08-c1-06-1 p
6. --- گرمی ساہ گھُٹݨ والی اے؛ کیوں جے \VUgras{کُتا} \textsuperscript{(15)} وی
ٹھنڈا ہوݨ آندا اے؛ اوہ \VUgras{نالے} \textsuperscript{(16)} دا تازہ پاݨی چٹدا اے
تے اک وچارے نِکے \VUgras{ڈڈو} \textsuperscript{(17)} نوں چونکا دیندا اے جیہڑا کنڈھے
دے گھاہ وچ دُبکیا ہویا سی۔ اوہ اوس صاف پاݨی وچ ٹپ پیندا اے جِتھے
\VUgras{کنول} \textsuperscript{(18)} کھِڑے نیں۔

%%K t08-c1-07-1 p
7. --- اک \VUgras{ممولا} \textsuperscript{(19)} اوس نِکے \VUgras{چشمے}
\textsuperscript{(20)} کول نہا رہیا اے جیہڑا دو چٹاناں وچوں پھُٹدا اے تے
\VUgras{کنڈیاں والیاں جھاڑیاں} \textsuperscript{(21)} تے \VUgras{بیریاں دیاں
جھاڑیاں} \textsuperscript{(22)} دے تھلے لُک جاندا اے۔

%%K t08-tit-3 sub
\VUsaut{5.0mm}
\VUpk{10.2pt}{40}{واڈھی دا کم۔}
\VUsaut{3.4mm}

%%K t08-c1-08-1 p
8. --- پنڈ دے نِکیاں لئی کݨک دی واڈھی وڈے پرچاوے دا ویلا اے۔ اوہ خوش
\VUgras{قلابازیاں} \textsuperscript{(24)} کھاندے نیں \VUgras{پرالی دے ڈھیراں}
\textsuperscript{(23)} اُتے، \VUgras{وڈھݨ والیاں} \textsuperscript{(26)} دی مدد کردے
نیں، تے جدوں اوہ \VUgras{کݨک دی پیلی} \textsuperscript{(27)} وڈھدے نیں اپݨیاں
\VUgras{دِرانیاں} \textsuperscript{(25)} نال، جیہڑیاں \VUgras{سان}
\textsuperscript{(30)} اُتے خوب تیز کیتیاں گئیاں نیں، تدوں نِکے کݨک کٹھی کر کے
\VUgras{پُولے} \textsuperscript{(31)} یا \VUgras{نِکیاں گٹھڑیاں} \textsuperscript{(32)}
بنھدے نیں۔

%%K t08-c1-09-1 p
9. --- کدے کدے اوہناں وچوں وڈیاں نوں \VUgras{درانتی} \textsuperscript{(12)} نال اوہ
\VUgras{سُنہری ڈنٹھل} \textsuperscript{(28)} وڈھݨ دِتے جاندے نیں جیہڑے داݨیاں نال
بھرے بھارے \VUgras{سِٹے} \textsuperscript{(29)} چُکی ہوئے نیں؛ تے نِکے،
\VUgras{ڈنگراں} \textsuperscript{(33)} دی رکھوالی کردے ہوئے، جیہڑے
\VUgras{میدان} \textsuperscript{(34)} وچ وڈے \VUgras{لِنڈن دے رُکھ}
\textsuperscript{(35)} دی چھاں وچ چردے نیں، اپݨی \VUgras{جالی} \textsuperscript{(36)}
وچ سوہݨیاں \VUgras{تِتلیاں} پھڑدے نیں۔

%%K t08-c1-09-2 p
کجھ تے \VUgras{گڈی} \textsuperscript{(37)} اُتے چڑھ کے اوہ پُولے جماندے نیں جیہڑے
اوہناں نوں \VUgras{پنجے} \textsuperscript{(38)} اُتے چُک کے دِتے جاندے نیں۔

%%K t08-c1-10-1 p
10. --- سارے \VUgras{میدان} \textsuperscript{(39)} وچ، جِتھوں \VUgras{چنڈولاں}
\textsuperscript{(41)} اُڈدیاں نیں، وڈی چہل پہل اے۔ سارے واڈھی وچ لگے نیں۔ پر
جِتھے غریب پرانے اوزار ای کم وچ لیاندے نیں — \VUgras{دِرانیاں، درانتیاں} تے
\VUgras{مُنگلی} \textsuperscript{(40)} —، اوتھے دھنی زمیندار اپݨی کݨک یا اپݨی
\VUgras{رائی} \textsuperscript{(43)} \VUgras{وڈھݨ دی مشین} \textsuperscript{(42)} نال
وڈھواندا اے۔ فیر، پُولیاں نوں کوٹھے تیکر ڈھوئے بغیر، اوتھے ای وڈے \VUgras{ڈھیر}
\textsuperscript{(44)} لا دِتے جاندے نیں۔

%%K t08-c1-11-1 p
11. --- فیر اک \VUgras{گاہݨ دی مشین} \textsuperscript{(47)} لیائی جاندی اے، جیہنوں
اک \VUgras{چلدا انجن} \textsuperscript{(45)} اک \VUgras{پٹے} \textsuperscript{(46)} دے
ذریعے چلاندا اے، تے ایس طرحاں داݨا \VUgras{پرالی} توں وکھ ہو جاندا اے، اوسے پیلی
وچ جِتھے اوہ بیجیا گیا سی۔

%%K t08-tit-4 sub
\VUsaut{5.0mm}
\VUpk{10.2pt}{40}{طوفان۔}
\VUsaut{3.4mm}

%%K t08-c1-12-1 p
12. --- واڈھی دے ویلے اکثر تِکھے \VUgras{طوفان} \textsuperscript{(53)} ٹُٹ پیندے
نیں۔ پہلاں دور توں \VUgras{بدلاں دی گرج} سُݨائی دیندی اے؛ اک تیز
\VUgras{لہندی ہوا} \textsuperscript{(54)} اُٹھدی اے تے \VUgras{جنگل}
\textsuperscript{(49)} دے سب توں وڈے رُکھاں نوں کراہندے ہوئے جھُکا دیندی اے۔ کالے
\VUgras{بدل} \textsuperscript{(56)} اسمان اُتے چڑھ آندے نیں؛ اوہناں نوں
\VUgras{بجلی} \textsuperscript{(55)} دیاں چُندھیاؤݨ والیاں ٹیڈھیاں لکیراں چیردیاں
نیں۔ \VUgras{مینہہ} تے کدے \VUgras{گڑے} پیݨ لگ پیندے نیں، تے وڈھݨ نوں تیار فصل
نوں وِچھا دیندے نیں۔

%%K t08-c1-13-1 p
13. --- اکثر بجلی کِسے عمارت نوں اگ لا دیندی اے۔ پچھلے ورھے، مثال لئی،
\VUgras{پون چکی} \textsuperscript{(50)} دی چھت سڑ کے سواہ ہو گئی تے اوہدے دو
\VUgras{پر} \textsuperscript{(51)} چُور چُور ہو گئے۔

%%K t08-c1-14-1 p
14. --- کجھ گھنٹیاں مگروں شانتی مُڑ آندی اے۔ اُفق اُتے اک چمکدی \VUgras{پینگھ}
\textsuperscript{(52)} وکھائی دیندی اے، تے قدرت اوہ جیوݨ مُڑ چُک لیندی اے جیہڑا
تھوڑی دیر لئی رُک گیا سی۔ جیہڑے پنڈ والے \VUgras{جھگیاں} \textsuperscript{(48)} وچ
آ لُکے سن، اوہ کم اُتے مُڑدے نیں تے ہُݨے ہُݨے ہوئے نقصان دی مرمت وچ ہمت نال جُٹ
جاندے نیں۔

%%K t08-tit-5 sub
\VUsaut{6.0mm}
\VUcentre{13.2pt}{90}{\textit{دوجا منظر۔}}
\VUsaut{3.6mm}

%%K t08-tit-6 sub
\VUpk{10.2pt}{40}{شکار۔ --- انگوراں دی واڈھی۔}
\VUsaut{1.4mm}
\VUpk{10.2pt}{40}{مچھی پھڑنا۔}
\VUsaut{4.0mm}

%%K t08-c2-01-1 p
1. --- \VUgras{اگست} دے چھیکر یا \VUgras{ستمبر} دے وچکار، علاقے مطابق، شکار دا
موسم کھلدا اے۔ سورج دیاں پہلیاں کِرناں \VUgras{کِرلیاں} \textsuperscript{(2)} نوں
اوہناں دیاں کھُڈاں توں باہر لیاندیاں وی نئیں پئی \VUgras{شکاری}
\textsuperscript{(5)} اپݨے \VUgras{کُتیاں} \textsuperscript{(4)} نال نکل پیندا اے۔

%%K t08-c2-02-1 p
2. --- اوہنے موٹے کپڑے دا اپݨا تکڑا \VUgras{شکاری سوٹ} \textsuperscript{(9)} پایا
اے، تے چمڑے دے \VUgras{پٹے}، جیہڑے اوہنوں اوس گِلے گھاہ وچ ٹُرن دیݨ گے جِتھے
\VUgras{ڈڈو} \textsuperscript{(1)} لُکے رہندے نیں؛ اوہنے اپݨی \VUgras{بندوق}
\textsuperscript{(6)} تے اپݨا \VUgras{شکار دا تھیلا} \textsuperscript{(10)} چُکیا اے،
تے ٹُر پیا اے۔

%%K t08-c2-03-1 p
3. --- شکار دی اُمنگ اوہنوں اک انگور دے باغ وچکار لے آندی اے جِتھے انگوراں دی
واڈھی پورے زور اُتے اے۔ انگور اُگاؤݨ والے دے نِکے، جیہڑے عام طور اُتے سڑک دے کنڈھے
بکریاں چاردے نیں، اج اوہناں نوں جِتھے جی کرے چرن دے رہے نیں، تے اک بُڈھے
\VUgras{بکرے} \textsuperscript{(3)} نوں کِلے نال بنھ کے — جیہڑا اپݨی آزادی دا فائدہ
چُک کے ضرور نس جاندا — اوہ وی اپݨے حصے دا مزہ لے رہے نیں۔ اک
\VUgras{انگوراں دی ٹوکری} \textsuperscript{(12)} توں اک گُچھا لیندا اے تے اوہدا
\VUgras{رَس} \textsuperscript{(14)} اک \VUgras{پیالے} \textsuperscript{(13)} وچ نچوڑ
کے مست (بغیر خمیر دی شراب) بݨاؤݨ وچ دل پرچاندا اے۔ دوجا اپݨے سر اُتے
\VUgras{پتّراں دا تاج} \textsuperscript{(15)} رکھ رہیا اے جیہڑا اوہنے ہُݨے گُندھیا
اے۔ اوہناں کول ڈھیٹ \VUgras{چڑیاں} \textsuperscript{(16)} انگور چوری کرن کولھو تیکر
چلیاں آندیاں نیں۔

%%K t08-c2-04-1 p
4. --- جدوں تیکر نِکے دل پرچاندے نیں، بندے کم کردے نیں۔
\VUgras{انگور توڑن والے} \textsuperscript{(18)} انگور \VUgras{پِٹھ دیاں ٹوکریاں}
\textsuperscript{(17)} وچ تہہ خانے تیکر لے جاندے نیں؛ اک \VUgras{پیپا ساز}
\textsuperscript{(19)} \VUgras{پیپیاں} \textsuperscript{(20)} دی مرمت کردا اے،
ویکھدا اے پئی \VUgras{پھٹے} \textsuperscript{(22)} تے \VUgras{ڈھکݨ}
\textsuperscript{(21)} ٹھیک نیں، تے ٹُٹے \VUgras{چھلے} \textsuperscript{(23)} بدلدا
اے۔ اوہ وڈے \VUgras{کُنڈ} \textsuperscript{(25)} وی اوسے نے بݨائے سن، جیہناں وچ خمیر
اُٹھاؤݨ نوں \VUgras{انگور} \textsuperscript{(26)} پائے جاندے نیں۔

%%K t08-c2-05-1 p
5. --- خمیر اُٹھ چُکݨ اُتے شراب کڈ لئی جاندی اے تے بچیا ہویا \VUgras{کولھو}
\textsuperscript{(28)} اُتے رکھیا جاندا اے؛ وڈے \VUgras{پیچ} \textsuperscript{(29)} نوں
ہولی ہولی کس کے رَس نچوڑیا جاندا اے، جیہڑا اک \VUgras{ٹب} \textsuperscript{(32)} وچ
وگدا اے تے ہور وی \VUgras{شراب} \textsuperscript{(31)} دیندا اے۔

%%K t08-tit-8 sub
\VUsaut{6.0mm}
\VUpk{10.2pt}{40}{بیئر تے سیب دی شراب۔}
\VUsaut{3.4mm}

%%K t08-c2-06-1 p
6. --- \VUgras{تہہ خانے} \textsuperscript{(27)} دی کندھ اُتے \VUgras{ہاپس}
\textsuperscript{(30)} دی اک ٹہݨی چڑھدی اے۔ ایتھے اوہ نِرا سجاوٹی بُوٹا اے، پر کجھ
دیساں وچ، جِتھے ویل بہت گھٹ ہوندی اے، \VUgras{بیئر} بݨاؤݨ لئی ہاپس اُگائے جاندے
نیں۔

%%K t08-c2-07-1 p
7. --- نِکا \VUgras{سیب دا رُکھ} \textsuperscript{(33)} سیباں نال لدیا اے، جیہناں
توں پکݨ اُتے ودھیا \VUgras{سیب دی شراب} بݨے گی۔ اپݨے \VUgras{پنجرے}
\textsuperscript{(34)} وچ \VUgras{کناری} \textsuperscript{(35)} تے
\VUgras{گولڈ فنچ} \textsuperscript{(36)} اوہناں چڑیاں نوں سڑ نال ویکھدے نیں جیہڑیاں
آزاد ٹپوسیاں مار رہیاں نیں۔

%%K t08-c2-08-1 p
8. --- جِتھوں تیکر نظر جاندی اے، ٹبیاں دی ڈھلوان اُتے، \VUgras{ویل دے کِلیاں}
\textsuperscript{(40)} دیاں قطاراں کھڑیاں نیں، جیہڑیاں شاندار \VUgras{ویل دے مُڈھاں}
\textsuperscript{(39)} نوں سنبھالی ہوئیاں نیں، \VUgras{گُچھیاں} نال لدیاں ہوئیاں۔

%%K t08-c2-08-2 p
سارا ورھا \VUgras{انگور اُگاؤݨ والے} \textsuperscript{(42)} نے اپݨے
\VUgras{انگور دے باغ} \textsuperscript{(38)} دی سنبھال وچ محنت کیتی اے، تے ہُݨ
اوہنوں چنگی فصل توں اوہدا انعام ملدا اے۔ \VUgras{توڑن والیاں}
\textsuperscript{(41)} اوہ کُنڈ بھردیاں نیں جیہڑے \VUgras{خچراں}
\textsuperscript{(43)} نال کھِچی گڈی اُتے رکھے نیں، تے سارے خوش نیں۔

%%K t08-tit-9 sub
\VUsaut{5.0mm}
\VUpk{10.2pt}{40}{مچھی پھڑنا۔}
\VUsaut{3.4mm}

%%K t08-c2-09-1 p
9. --- ایہو گل \VUgras{بنسی والیاں} \textsuperscript{(44)} دی اے، جیہڑے پوہ پھٹدے
ای اپݨیاں \VUgras{بنسیاں} \textsuperscript{(45)} نال پاݨی دے کنڈھے آ جمے نیں۔

%%K t08-c2-09-2 p
\VUgras{مچھیاں} \textsuperscript{(47)} کنڈا پھڑ لیندیاں نیں جُوں ای اوہ اپݨی
\VUgras{ڈور} \textsuperscript{(46)} پاݨی وچ پاندے نیں، تے اوہناں نوں \VUgras{چوگا}
بدلݨ دا ویلا وی نئیں ملدا پئی ڈور دے سرے اُتے اک نواں شکار تڑپ رہیا ہوندا اے۔

%%K t08-c2-09-3 p
فیر وی \VUgras{جال والا ماہی گیر} \textsuperscript{(48)}، جیہڑا اپݨی \VUgras{بیڑی}
\textsuperscript{(49)} توں \VUgras{دریا} \textsuperscript{(50)} دے وچکار جال سُٹدا اے،
کِتے ودھ پھڑدا اے۔

%%K t08-c2-10-1 p
10. --- پتجھڑ چنگی سیر دا موسم اے : پیدل، \VUgras{گھوڑے} \textsuperscript{(52)}
اُتے، \VUgras{سائیکل} \textsuperscript{(53)} نال، \VUgras{موٹر} \textsuperscript{(54)}
نال۔ \VUgras{سڑکاں} \textsuperscript{(51)} صاف نیں تے لوک چھیکرلے چنگے دناں دا پورا
فائدہ چُکدے نیں، کیوں جے \VUgras{ابابیلاں} \textsuperscript{(56)} پہلاں ای چھتاں اُتے
جُڑن لگ پئیاں نیں، نیمے دیساں نوں اُڈ جاوݨ لئی۔ بُڈھے \VUgras{اخروٹ دے رُکھ}
\textsuperscript{(57)} دے پتّر پیلے پے رہے نیں، \VUgras{اخروٹ} ڈِگ رہے نیں، تے اک
\VUgras{جھُنڈ} \textsuperscript{(58)} وچ کھڑے اُچے \VUgras{رُکھ}، جیہڑا
\VUgras{اُچی زمین} \textsuperscript{(59)} اُتے اے، چھیتی ای ننگے ہو جاوݨ گے۔

%%K t08-c2-11-1 p
11. --- \VUgras{سورج} \textsuperscript{(60)} دیر نال چڑھدا اے، دن نِکے ہوندے جاندے
نیں، \VUgras{سویر} اک چنگی خاصی ٹھنڈ محسوس ہوݨ لگ پیندی اے، تے \VUgras{چاہاں}
\textsuperscript{(61)} دے جھُنڈ ہُݨے توں ساڈا بے رحم موسم چھڈ کے جا رہے نیں۔
\VUsaut{6.0mm}
\VUfilet{22mm}
